% GENERATED FILE. Edit canonical lessons, then run npm run generate:books.
% canonical-source-hash: fnv1a64:e19668fa289b3ab5
% canonical-lessons: HI-C501-nata, HI-C501-byahna, HI-C501-talaq, HI-C501-nibhana, HI-C501-sulah

\chapter{Getting Along}
\label{ch:hi-getting-along}
\hlchaptermodality{501}

\begin{chapteropening}
\textbf{By the end of this chapter:} \emph{I can say how family members get on.}

Complete the last lesson of chapter 501: I can say how family members get on.
\end{chapteropening}

\section[\texorpdfstring{\hi{नाता} (nātā)}{नाता (nātā)}]{\hi{नाता} (nātā) — a bond; a tie of kinship}

\label{lesson:HI-C501-nata}



\begin{quote}
Before the new one: say the Hindi for the groom's wedding procession, then the Hindi for a farewell.
\end{quote}

\subsection*{You'll want to know: \hi{नाता}}
\textbf{\hi{नाता}} — \emph{nātā} — \textquotedblleft{}a bond; a tie of kinship\textquotedblright{}.

\begin{quote}
\hi{मीरा} \hi{का} \hi{अपने} \hi{गाँव} \hi{से} \hi{गहरा} \hi{नाता} \hi{है।}
\end{quote}

\begin{quote}
\emph{Mīrā kā apne gā̃v se gahrā nātā hai.}
\end{quote}

\begin{quote}
Mira has a deep bond with her village.
\end{quote}

Masculine.

\begin{cousinweb}
\textbf{\hi{नाता}} (\emph{nātā}) is a tadbhava: CDIAL lists Hindi \emph{nātā}, \textquotedblleft{}kinship, alliance\textquotedblright{}, under \emph{*jñātra}, a form it reconstructs, close to Sanskrit \emph{jñāti}, \textquotedblleft{}a kinsman\textquotedblright{} (CDIAL 5279).
\end{cousinweb}

\begin{grammarlens}[title={What you've built}]
One more word for describing what you lived and felt.
\end{grammarlens}

\subsection*{Your turn}
\begin{itemize}
\raggedright
  \item \emph{Say it:} \emph{nātā}
  \item \emph{Say it:} \emph{nātā}, once more
  \item \emph{Say it:} \emph{Mīrā kā apne gā̃v se gahrā nātā hai.} — Mira has a deep bond with her village.
  \item \emph{Recall:} say the Hindi for the groom's wedding procession, then the Hindi for a farewell, then say \emph{nātā} again
\end{itemize}

\begin{tcolorbox}[breakable,colback=teal!4,colframe=teal!35!black,title={Before you move on}]
What does \hi{नाता} mean? (\textquotedblleft{}A bond; a tie of kinship\textquotedblright{}.) Say it once more.
\end{tcolorbox}
\section[\texorpdfstring{\hi{ब्याहना} (byāhnā)}{ब्याहना (byāhnā)}]{\hi{ब्याहना} (byāhnā) — to marry; to marry off}

\label{lesson:HI-C501-byahna}



\begin{quote}
Before the new one: say the Hindi for a farewell, then the Hindi for a bond.
\end{quote}

\subsection*{You'll want to know: \hi{ब्याहना}}
\textbf{\hi{ब्याहना}} — \emph{byāhnā} — \textquotedblleft{}to marry; to marry off\textquotedblright{}.

\begin{quote}
\hi{पिता} \hi{ने} \hi{बेटी} \hi{को} \hi{अच्छे} \hi{परिवार} \hi{में} \hi{ब्याहा।}
\end{quote}

\begin{quote}
\emph{pitā ne beṭī ko acche parivār meṁ byāhā.}
\end{quote}

\begin{quote}
The father married his daughter into a good family.
\end{quote}

Past, with \hi{ने} (he or she …): \hi{उसने} \hi{बेटी} \hi{को} \hi{ब्याहा}, \emph{usne beṭī ko byāhā}.

\begin{cousinweb}
\textbf{\hi{ब्याहना}} (\emph{byāhnā}) is a tadbhava: CDIAL lists Hindi \emph{byāhnā}, \textquotedblleft{}to marry\textquotedblright{}, under Sanskrit \emph{vivāhayati} (CDIAL 11923), the verb beside \emph{vivāha}, \textquotedblleft{}marriage\textquotedblright{}.
\end{cousinweb}

\begin{grammarlens}[title={What you've built}]
One more word for describing what you lived and felt.
\end{grammarlens}

\subsection*{Your turn}
\begin{itemize}
\raggedright
  \item \emph{Say it:} \emph{byāhnā}
  \item \emph{Say it:} \emph{byāhnā}, once more
  \item \emph{Say it:} \emph{pitā ne beṭī ko acche parivār meṁ byāhā.} — The father married his daughter into a good family.
  \item \emph{Recall:} say the Hindi for a farewell, then the Hindi for a bond, then say \emph{byāhnā} again
\end{itemize}

\begin{tcolorbox}[breakable,colback=teal!4,colframe=teal!35!black,title={Before you move on}]
What does \hi{ब्याहना} mean? (\textquotedblleft{}To marry; to marry off\textquotedblright{}.) Say it once more.
\end{tcolorbox}
\section[\texorpdfstring{\hi{तलाक़} (talāq)}{तलाक़ (talāq)}]{\hi{तलाक़} (talāq) — a divorce}

\label{lesson:HI-C501-talaq}



\begin{quote}
Before the new one: say the Hindi for a bond, then the Hindi for to marry.
\end{quote}

\subsection*{You'll want to know: \hi{तलाक़}}
\textbf{\hi{तलाक़}} — \emph{talāq} — \textquotedblleft{}a divorce\textquotedblright{}.

\begin{quote}
\hi{दो} \hi{साल} \hi{बाद} \hi{पति} \hi{और} \hi{पत्नी} \hi{का} \hi{तलाक़} \hi{हो} \hi{गया।}
\end{quote}

\begin{quote}
\emph{do sāl bād pati aur patnī kā talāq ho gayā.}
\end{quote}

\begin{quote}
Two years later the husband and wife got divorced.
\end{quote}

Masculine.

\begin{cousinweb}
\textbf{\hi{तलाक़}} (\emph{talāq}) is usually traced to Arabic \emph{ṭalāq}, \textquotedblleft{}release, divorce\textquotedblright{}, which reached Hindi through Persian.
\end{cousinweb}

\begin{grammarlens}[title={What you've built}]
One more word for describing what you lived and felt.
\end{grammarlens}

\subsection*{Your turn}
\begin{itemize}
\raggedright
  \item \emph{Say it:} \emph{talāq}
  \item \emph{Say it:} \emph{talāq}, once more
  \item \emph{Say it:} \emph{do sāl bād pati aur patnī kā talāq ho gayā.} — Two years later the husband and wife got divorced.
  \item \emph{Recall:} say the Hindi for a bond, then the Hindi for to marry, then say \emph{talāq} again
\end{itemize}

\begin{tcolorbox}[breakable,colback=teal!4,colframe=teal!35!black,title={Before you move on}]
What does \hi{तलाक़} mean? (\textquotedblleft{}A divorce\textquotedblright{}.) Say it once more.
\end{tcolorbox}
\section[\texorpdfstring{\hi{निभाना} (nibhānā)}{निभाना (nibhānā)}]{\hi{निभाना} (nibhānā) — to keep up; to stand by}

\label{lesson:HI-C501-nibhana}



\begin{quote}
Before the new one: say the Hindi for to marry, then the Hindi for a divorce.
\end{quote}

\subsection*{You'll want to know: \hi{निभाना}}
\textbf{\hi{निभाना}} — \emph{nibhānā} — \textquotedblleft{}to keep up; to stand by\textquotedblright{}.

\begin{quote}
\hi{अरुण} \hi{ने} \hi{अपना} \hi{वादा} \hi{निभाया।}
\end{quote}

\begin{quote}
\emph{Aruṇ ne apnā vādā nibhāyā.}
\end{quote}

\begin{quote}
Arun kept his promise.
\end{quote}

Past, with \hi{ने} (he or she …): \hi{उसने} \hi{वादा} \hi{निभाया}, \emph{usne vādā nibhāyā}.

\begin{cousinweb}
\textbf{\hi{निभाना}} (\emph{nibhānā}) is the transitive beside \emph{nibhnā}, \textquotedblleft{}to last, succeed\textquotedblright{}: CDIAL lists \emph{nibhnā} and \emph{nibāhnā}, \textquotedblleft{}to carry on\textquotedblright{}, under Sanskrit \emph{nírvahati}, \textquotedblleft{}carries through\textquotedblright{} (CDIAL 7397). One keeps up a promise, a friendship or a duty.
\end{cousinweb}

\begin{grammarlens}[title={What you've built}]
One more word for describing what you lived and felt.
\end{grammarlens}

\subsection*{Your turn}
\begin{itemize}
\raggedright
  \item \emph{Say it:} \emph{nibhānā}
  \item \emph{Say it:} \emph{nibhānā}, once more
  \item \emph{Say it:} \emph{Aruṇ ne apnā vādā nibhāyā.} — Arun kept his promise.
  \item \emph{Recall:} say the Hindi for to marry, then the Hindi for a divorce, then say \emph{nibhānā} again
\end{itemize}

\begin{tcolorbox}[breakable,colback=teal!4,colframe=teal!35!black,title={Before you move on}]
What does \hi{निभाना} mean? (\textquotedblleft{}To keep up; to stand by\textquotedblright{}.) Say it once more.
\end{tcolorbox}
\section[\texorpdfstring{\hi{सुलह} (sulah)}{सुलह (sulah)}]{\hi{सुलह} (sulah) — a reconciliation; making peace}

\label{lesson:HI-C501-sulah}



\begin{quote}
Before the new one: say the Hindi for a divorce, then the Hindi for to keep up.
\end{quote}

\subsection*{You'll want to know: \hi{सुलह}}
\textbf{\hi{सुलह}} — \emph{sulah} — \textquotedblleft{}a reconciliation; making peace\textquotedblright{}.

\begin{quote}
\hi{बड़ों} \hi{ने} \hi{दोनों} \hi{परिवारों} \hi{में} \hi{सुलह} \hi{कराई।}
\end{quote}

\begin{quote}
\emph{baṛoṁ ne donõ parivāroṁ meṁ sulah karāī.}
\end{quote}

\begin{quote}
The elders brought the two families back together.
\end{quote}

Feminine.

\begin{cousinweb}
\textbf{\hi{सुलह}} (\emph{sulah}) is usually traced to Arabic \emph{ṣulḥ}, \textquotedblleft{}peace, reconciliation\textquotedblright{}, which reached Hindi through Persian.
\end{cousinweb}

\begin{grammarlens}[title={What you've built}]
One more word for describing what you lived and felt.
\end{grammarlens}

\subsection*{Your turn}
\begin{itemize}
\raggedright
  \item \emph{Say it:} \emph{sulah}
  \item \emph{Say it:} \emph{sulah}, once more
  \item \emph{Say it:} \emph{baṛoṁ ne donõ parivāroṁ meṁ sulah karāī.} — The elders brought the two families back together.
  \item \emph{Recall:} say the Hindi for a divorce, then the Hindi for to keep up, then say \emph{sulah} again
\end{itemize}

\begin{tcolorbox}[breakable,colback=teal!4,colframe=teal!35!black,title={Before you move on}]
What does \hi{सुलह} mean? (\textquotedblleft{}A reconciliation; making peace\textquotedblright{}.) Say it once more.
\end{tcolorbox}
